%==============================
\chapter{Introduction}\label{cap:introducao}
%==============================

The widespread adoption of multicore and many-core processor architectures has established parallel programming as a foundational skill in high-performance computing (HPC). 
Modern commodity processors routinely provide six, eight, or more cores, and larger systems scale to dozens or hundreds of cores in a single node. 
Exploiting this hardware requires programs that can distribute work across multiple threads executing concurrently, yet sequential programming remains the dominant paradigm for scientific and engineering software. 
The gap between available hardware parallelism and the sequential nature of most existing code represents a persistent challenge for the HPC community.

Converting a sequential program to a parallel one requires more than inserting directives into the source code. 
The programmer must identify independent computations, reason about data sharing and race conditions, select appropriate synchronization primitives, and ensure that the parallel version produces results identical to those of the original sequential code. 
These tasks are error-prone even for experienced practitioners. 
Mistakes such as incorrect scoping of loop variables, missing barriers between dependent phases, or deadlocks in synchronization can produce code that compiles and sometimes produces plausible output, yet is functionally incorrect or consistently slower than the sequential baseline.

Large Language Models (LLMs) have recently demonstrated substantial code generation capabilities. 
Models such as ChatGPT, Gemini, Claude, Meta AI, and DeepSeek are trained on large corpora of text and source code.
They are capable of generating, explaining, and transforming programs across a variety of languages and programming paradigms. 
A natural question arising from these capabilities is whether LLMs can automate one of the most error-prone and expertise-intensive steps in software development: converting sequential C code into parallel code using OpenMP, the most widely used programming interface for shared-memory parallelism.

This work investigates that question through a controlled empirical evaluation. 
Five LLMs were each asked to parallelize six sequential C programs drawn from the PAMPAR benchmark \citeonline{GARCIA:PAMPAR:2020}, which provides validated sequential implementations together with expert-written reference parallel versions. 
The generated programs were evaluated for compilation success, output correctness, and execution speedup relative to the sequential baseline, using 11 threads on a shared-memory machine. 
Hardware performance counters were also collected to support a deeper analysis of the observed performance differences.

%==============================
\section{Objectives}
\label{sec:objectives}

The primary objective of this work is to evaluate the capability of current LLMs to convert sequential C code to parallel code using OpenMP. 
Specifically, this work aims to:

\begin{alineas}
\item Assess whether the parallel code generated by LLMs compiles without errors and produces output that matches the sequential original;
\item Measure the speedup achieved by the generated parallel code at 11 threads, using the sequential baseline and the PAMPAR reference parallel implementation as reference points;
\item Characterize the most common failure modes and limitations present in the generated code, including incorrect parallelization strategies that compile and produce correct output but achieve no performance gain;
\item Compare the five evaluated LLMs in terms of correctness and performance of the generated parallel code.
\end{alineas}

%==============================
\section{Contributions}
\label{sec:contributions}

The main contributions of this work are:

\begin{alineas}
\item A controlled evaluation of five publicly accessible LLMs across six benchmark programs that span a range of parallelization difficulty, from embarrassingly parallel reductions to problems with wavefront data dependencies;
\item A systematic correctness and performance analysis of all generated programs, grounded in the PAMPAR reference implementations as an expert-written performance target;
\item An identification and characterization of recurring failure patterns in LLM-generated parallel code, providing concrete examples of how current models can produce compilable, correct code that nonetheless fails to achieve meaningful parallelism.
\end{alineas}

%==============================
\section{Document Structure}
\label{sec:structure}

The remainder of this work is organized as follows. 
Chapter~\ref{chapter:trabalhos} discusses related work on LLM-based code generation, focusing on studies that evaluate LLMs for parallel and HPC code. 
Chapter~\ref{cap:background} presents the theoretical background covering parallel programming, OpenMP, performance metrics, and the PAMPAR benchmark. 
Chapter~\ref{cha:methodology} describes the research questions, the selected programs and models, the prompting strategy, and the evaluation procedure. 
Chapter~\ref{cap:resultado} presents and analyzes the experimental results organized around the four research questions. 
Chapter~\ref{cap:conclusao} summarizes the findings and outlines directions for future research.